In this part you will use the sky-map found in the envelope. The map
represents the sky in Suceava (latitude $47^\circ 39'$ North, longitude
$26^\circ 15'$ East) at 19:00 UT on the day of the test. The observer who
made the sky-map was at a very high altitude above Suceava; the Zenith point
is in the centre of the chart. Use a pencil for marking and drawing lines on
the sky-map. Solve questions 1 to 4 on one copy of the map and questions 5
to 8 on the second copy of the map. You have 30 minutes to finish this part.

\begin{description}
\item[(1)] (2p) Draw on the map the horizon for an observer located on the
  ground in Suceava.
\item[(2)] (8p) Draw the celestial equator, the ecliptic, the galactic
  equator and the local meridian on the map with continuous lines.
\item[(3)] (9p) Mark the cardinal points (as N for north, E for east, S for
  south and W for west). Mark all the visible planets (except Uranus and
  Neptune) of the Solar System on the map and number them as 1, 2, \ldots,
  6 in the order of increasing orbital radius (skip number 3 for the Earth).
  Note that planets are not currently shown on the map.
\item[(4)] (4p) Identify and mark the four brightest stars in the visual
  band above the horizon line. Number the stars starting from 1 --- the
  brightest, and continue with the fainter ones till number 4 for the
  faintest. Fill in a table with the Bayer names of the four identified
  stars.
\item[(5)] (6p) Draw on the map approximate figures of any 15 constellations
  which lie completely above the horizon. Each constellation you mark should
  be identified on the map with the IAU abbreviation, using the provided
  table of constellations.
\item[(6)] (5p) Mark on the map the positions of the following objects:
  \begin{description}
  \item[(a)] the Messier objects M31, M27, M13;
  \item[(b)] $\gamma$ Cygni, $\beta$ Ursae Minoris.
  \end{description}
\item[(7)] (10p) Estimate the sidereal time of the map; write the value in
  the box.
\item[(8)] (6p) Estimate the equatorial coordinates (right ascension
  $\alpha$ and declination $\delta$) of the star Altair ($\alpha$ Aquilae).
\end{description}

% FIGURE NEEDED: zenith-centred sky map of Suceava at 19:00 UT on the day of
% the test (two copies), on which all marking tasks are carried out
